\documentclass[11pt]{article}
\usepackage{amsmath,amssymb,amsthm}
\usepackage{algorithm}
\usepackage[noend]{algpseudocode}
\usepackage{fullpage}
\usepackage{enumitem}
\newtheorem{theorem}{Theorem}
\newtheorem{lemma}{Lemma}
\newtheorem{assumption}{Assumption}
\newcommand{\sign}{\operatorname{sign}}
\newcommand{\E}{\mathbb{E}}
\newcommand{\R}{\mathbb{R}}

\begin{document}

\section*{SignSGD with Momentum}

\section*{Mathematical Formulation}
Let $f:\R^{d}\rightarrow\R$ be a (possibly non‑convex) differentiable loss.
At iteration $t\in\{0,1,\dots\}$ we maintain a momentum vector $m_t\in\R^{d}$
and a parameter vector $x_t\in\R^{d}$.  
Given a stochastic gradient $g_t$ (e.g. a minibatch gradient) we define  

\[
\boxed{
\begin{aligned}
m_{t+1} &= \beta\, m_{t} + (1-\beta)\, g_{t}, \qquad \beta\in[0,1),\\[4pt]
x_{t+1} &= x_{t} - \alpha\, \sign\!\big(m_{t+1}\big),
\end{aligned}}
\]
where $\alpha>0$ is a stepsize and $\sign(v)$ is applied element‑wise,
$\sign(v)_i = \begin{cases}+1 & v_i>0\\ -1 & v_i<0\\ 0 & v_i=0.\end{cases}$

\section*{Pseudocode}
\begin{algorithm}[H]
\caption{SignSGD with Momentum}
\begin{algorithmic}[1]
\State \textbf{Input:} stepsize $\alpha>0$, momentum $\beta\in[0,1)$,
               max iterations $T$
\State Initialize $x_0\in\R^{d}$, $m_0 = 0$
\For{$t=0$ \textbf{to} $T-1$}
    \State Sample a minibatch $\mathcal{B}_t$ and compute stochastic gradient 
          $g_t = \nabla f(x_t; \mathcal{B}_t)$
    \State $m_{t+1} \gets \beta\, m_t + (1-\beta)\, g_t$ \hfill\Comment{momentum update}
    \State $x_{t+1} \gets x_t - \alpha\, \sign(m_{t+1})$ \hfill\Comment{signed step}
\EndFor
\State \textbf{Output:} $x_T$
\end{algorithmic}
\end{algorithm}

\section*{Dry Run Example}
Consider $f(w) = (w-3)^2$ (so $d=1$).  
The true gradient is $\nabla f(w)=2(w-3)$.  Use $\alpha=0.1$, $\beta=0.9$, and a deterministic “stochastic’’ gradient $g_t=\nabla f(w_t)$.

\begin{center}
\begin{tabular}{c|c|c|c|c}
$t$ & $w_t$ & $g_t$ & $m_{t+1}$ & $w_{t+1}=w_t-\alpha\,\sign(m_{t+1})$\\ \hline
0 & $0$   & $-6$ & $0.9\cdot0+(0.1)(-6)=-0.6$ & $0-0.1\cdot(-1)=0.1$\\
1 & $0.1$ & $-5.8$ & $0.9(-0.6)+0.1(-5.8)=-1.04$ & $0.1-0.1\cdot(-1)=0.2$\\
2 & $0.2$ & $-5.6$ & $0.9(-1.04)+0.1(-5.6)=-1.456$ & $0.2-0.1\cdot(-1)=0.3$\\
\end{tabular}
\end{center}
After three iterations the objective values are
$f(0)=9$, $f(0.1)=8.41$, $f(0.2)=7.84$, $f(0.3)=7.29$, showing descent.

\newpage
\section*{Assumptions}
\begin{assumption}[Smoothness]\label{ass:smooth}
$f$ has $L$‑Lipschitz gradients, i.e.
\[
\|\nabla f(x)-\nabla f(y)\|_2\le L\|x-y\|_2,\qquad\forall x,y\in\R^{d}.
\]
\end{assumption}

\begin{assumption}[Unbiased Stochastic Gradient]\label{ass:unbiased}
For each $t$, the stochastic gradient $g_t$ satisfies
\[
\E[g_t\mid x_t]=\nabla f(x_t),\qquad
\E\big[\|g_t-\nabla f(x_t)\|_2^{2}\mid x_t\big]\le\sigma^{2},
\]
for a constant $\sigma\ge0$.
\end{assumption}

\begin{assumption}[Bounded Gradient Coordinates]\label{ass:bounded}
There exists $G>0$ such that for all $t$,
\[
\|\,\nabla f(x_t)\|_{\infty}\le G\quad\text{and}\quad
\|g_t\|_{\infty}\le G\;\;\text{a.s.}
\]
\end{assumption}

\begin{assumption}[Stepsize]\label{ass:stepsize}
The stepsize $\alpha$ is constant and satisfies
\[
0<\alpha\le\frac{1}{L\sqrt{d}}.
\]
\end{assumption}

\section*{Main Convergence Theorem}
\begin{theorem}\label{thm:main}
Assume \ref{ass:smooth}--\ref{ass:stepsize} hold. Let $\{x_t\}$ be generated by SignSGD with Momentum.
Define the averaged gradient norm
\[
\bar{G}_T:=\frac{1}{T}\sum_{t=0}^{T-1}\E\big[\|\nabla f(x_t)\|_2\big].
\]
Then for any $T\ge1$,
\[
\bar{G}_T\;\le\;
\frac{2\big(f(x_0)-f^{\star}\big)}{\alpha T}
   \;+\;
\frac{L\alpha d}{2}
   \;+\;
\frac{2(1-\beta)G}{\sqrt{d}}.
\]
Consequently, choosing $\alpha =\Theta\!\big(\frac{1}{\sqrt{T}}\big)$ yields
\[
\bar{G}_T = O\!\left(\frac{1}{\sqrt{T}}\right).
\]
\end{theorem}

\section*{Theoretical Properties}
\begin{enumerate}[leftmargin=*]
\item \textbf{Stability.} The bound contains a term $L\alpha d/2$ that grows linearly with the dimension $d$, reflecting the fact that the signed step has fixed $\ell_2$‑norm $\sqrt{d}$. Choosing $\alpha\le 1/(L\sqrt{d})$ guarantees that this term does not dominate the $1/(\alpha T)$ term.
\item \textbf{Momentum Effect.} The additive term $\frac{2(1-\beta)G}{\sqrt{d}}$ vanishes as $\beta\to1$, showing that a high momentum reduces the bias introduced by the sign operator.
\item \textbf{Comparison to Vanilla SGD.} For SGD under the same assumptions one obtains 
$\bar{G}_T = O(1/\sqrt{T})$ with a constant that does not depend on $d$. Hence SignSGD sacrifices a factor $\sqrt{d}$ in the constant but gains communication efficiency because only one bit per coordinate is transmitted.
\end{enumerate}

\section*{Discussion}
The theorem demonstrates that, despite the extreme quantisation of the gradient to its sign, the algorithm still converges to a stationary point at the same $O(1/\sqrt{T})$ rate as full‑precision SGD, provided the stepsize is scaled as $1/\sqrt{d}$. Momentum mitigates the bias introduced by sign quantisation and improves the constant in the bound.

\newpage
\section*{Preliminaries (Definitions and Lemmas)}
\begin{lemma}[Signed Descent Lemma]\label{lem:signed-descent}
For any vectors $a,b\in\R^{d}$,
\[
a^{\top}\sign(b) \;\ge\; \frac{1}{\sqrt{d}}\|a\|_{2}\;-\;\|a-b\|_{2}.
\]
\end{lemma}
\begin{proof}
Since each coordinate of $\sign(b)$ belongs to $\{-1,0,1\}$,
\[
a^{\top}\sign(b)=\sum_{i=1}^{d}a_i\sign(b_i)
\ge\sum_{i=1}^{d}|a_i|-\sum_{i=1}^{d}|a_i-\sign(b_i)|
\ge\|a\|_{1}-\|a-\sign(b)\|_{1}.
\]
Using $\|a\|_{1}\ge\|a\|_{2}$ and $\|u\|_{1}\le\sqrt{d}\|u\|_{2}$ we obtain
\[
a^{\top}\sign(b)\ge\|a\|_{2}-\sqrt{d}\|a-\sign(b)\|_{2}.
\]
Finally note that $\|\sign(b)-b\|_{2}\le\|a-b\|_{2}+\|a-\sign(b)\|_{2}$, which
implies $\|a-\sign(b)\|_{2}\le\|a-b\|_{2}$. Substituting gives the claimed inequality. \qedhere
\end{proof}

\begin{lemma}[Momentum Bias]\label{lem:momentum-bias}
Define $m_{t+1}= \beta m_t+(1-\beta)g_t$ with $m_0=0$.
Then
\[
\E[m_{t+1}\mid x_t]= (1-\beta)\sum_{k=0}^{t}\beta^{t-k}\,\nabla f(x_k).
\]
Consequently,
\[
\big\|\E[m_{t+1}\mid x_t]-\nabla f(x_t)\big\|_{2}
\le (1-\beta)\,G\sum_{k=0}^{t-1}\beta^{t-1-k}
\le\frac{(1-\beta)G}{1-\beta}= (1-\beta)G .
\]
\end{lemma}
\begin{proof}
Unrolling the recursion yields
\[
m_{t+1}= \beta^{t+1}m_0 + (1-\beta)\sum_{k=0}^{t}\beta^{t-k} g_k.
\]
Since $m_0=0$ and $\E[g_k\mid x_k]=\nabla f(x_k)$, taking conditional expectation
conditioned on $x_t$ (which also conditions on all previous iterates) gives the first identity.
The second inequality follows from the boundedness of each gradient coordinate
(Assumption~\ref{ass:bounded}) and the geometric series formula. \qedhere
\end{proof}

\newpage
\section*{Main Proof (Step‑by‑Step Derivation)}
We prove Theorem~\ref{thm:main}.  All expectations are taken over the randomness of the stochastic gradients up to iteration $t$.

\begin{enumerate}
\item[Step 1.] \textbf{Smoothness inequality.}  
By Assumption~\ref{ass:smooth},
\[
f(x_{t+1})\le f(x_t)+\nabla f(x_t)^{\top}(x_{t+1}-x_t)
                +\frac{L}{2}\|x_{t+1}-x_t\|_2^{2}. \tag{1}
\]

\item[Step 2.] \textbf{Plug the update.}  
Using $x_{t+1}=x_t-\alpha\sign(m_{t+1})$,
\[
x_{t+1}-x_t = -\alpha\sign(m_{t+1}),\qquad
\|x_{t+1}-x_t\|_2^{2}= \alpha^{2}\|\sign(m_{t+1})\|_2^{2}= \alpha^{2}d.
\] 
Insert into (1):
\[
f(x_{t+1})\le f(x_t)-\alpha\,\nabla f(x_t)^{\top}\sign(m_{t+1})
            +\frac{L\alpha^{2}d}{2}. \tag{2}
\]

\item[Step 3.] \textbf{Take conditional expectation.}  
Condition on the sigma‑algebra $\mathcal{F}_t$ generated by $\{x_0,\dots,x_t,m_t\}$.
Denote $\E_t[\cdot]=\E[\cdot\mid\mathcal{F}_t]$.  Taking $\E_t$ on both sides of (2) gives
\[
\E_t[f(x_{t+1})]\le f(x_t)-\alpha\,\E_t\!\big[\nabla f(x_t)^{\top}\sign(m_{t+1})\big]
            +\frac{L\alpha^{2}d}{2}. \tag{3}
\]

\item[Step 4.] \textbf{Lower bound the inner product.}  
Apply Lemma~\ref{lem:signed-descent} with $a=\nabla f(x_t)$ and $b=m_{t+1}$:
\[
\nabla f(x_t)^{\top}\sign(m_{t+1})
   \ge \frac{1}{\sqrt{d}}\|\nabla f(x_t)\|_{2}
        -\|\nabla f(x_t)-m_{t+1}\|_{2}. \tag{4}
\]

\item[Step 5.] \textbf{Control the deviation $\|\nabla f-m\|$.}  
Decompose
\[
\|\nabla f(x_t)-m_{t+1}\|_{2}
\le \underbrace{\|\nabla f(x_t)-\E_t[m_{t+1}]\|_{2}}_{\text{bias}}
   + \underbrace{\|m_{t+1}-\E_t[m_{t+1}]\|_{2}}_{\text{variance}}.
\] 
Lemma~\ref{lem:momentum-bias} yields the bias bound
\[
\|\nabla f(x_t)-\E_t[m_{t+1}]\|_{2}\le (1-\beta)G. \tag{5}
\]
For the variance term we use Jensen’s inequality and the bounded variance of $g_t$:
\[
\E_t\big[\|m_{t+1}-\E_t[m_{t+1}]\|_{2}\big]
      \le \E_t\big[\|m_{t+1}-\E_t[m_{t+1}]\|_{2}^{2}\big]^{1/2}
      \le \frac{(1-\beta)\sigma}{\sqrt{1-\beta^{2}}}
      \le (1-\beta) \sigma . \tag{6}
\]
(Here we used the explicit expression of $m_{t+1}$ as a weighted sum of the $g_k$'s.)

\item[Step 6.] \textbf{Combine (4)–(6).}  
Taking expectation $\E_t$ of (4) and employing (5)–(6) gives
\[
\E_t\!\big[\nabla f(x_t)^{\top}\sign(m_{t+1})\big]
   \ge \frac{1}{\sqrt{d}}\|\nabla f(x_t)\|_{2}
        -(1-\beta)G-(1-\beta)\sigma . \tag{7}
\]

\item[Step 7.] \textbf{Insert (7) into (3).}  
\[
\E_t[f(x_{t+1})]\le f(x_t)
   -\frac{\alpha}{\sqrt{d}}\|\nabla f(x_t)\|_{2}
   +\alpha(1-\beta)(G+\sigma)
   +\frac{L\alpha^{2}d}{2}. \tag{8}
\]

\item[Step 8.] \textbf{Take total expectation and rearrange.}  
Denote $\Delta_t:=\E[f(x_t)]-f^{\star}$ where $f^{\star}= \inf_x f(x)$.  Taking full expectation of (8) and moving the gradient term to the left yields
\[
\frac{\alpha}{\sqrt{d}}\E\!\big[\|\nabla f(x_t)\|_{2}\big]
   \le \Delta_t-\Delta_{t+1}
      +\alpha(1-\beta)(G+\sigma)
      +\frac{L\alpha^{2}d}{2}. \tag{9}
\]

\item[Step 9.] \textbf{Sum over $t=0,\dots,T-1$.}  
Summing (9) telescopes the $\Delta$ terms:
\[
\frac{\alpha}{\sqrt{d}}\sum_{t=0}^{T-1}\E\!\big[\|\nabla f(x_t)\|_{2}\big]
   \le \Delta_0-\Delta_T
      +T\alpha(1-\beta)(G+\sigma)
      +\frac{L\alpha^{2}d}{2}T . \tag{10}
\]
Since $\Delta_T\ge0$, we drop it and divide by $T\alpha/\sqrt{d}$:
\[
\frac{1}{T}\sum_{t=0}^{T-1}\E\!\big[\|\nabla f(x_t)\|_{2}\big]
   \le \frac{\sqrt{d}\,\Delta_0}{\alpha T}
      + (1-\beta)(G+\sigma)
      + \frac{L\alpha d^{3/2}}{2}. \tag{11}
\]

\item[Step 10.] \textbf{Simplify constants.}  
Define $C_1:=2\Delta_0$, $C_2:=L/2$, and absorb $\sigma$ into $G$ (both are bounded by a universal constant under Assumption~\ref{ass:bounded}).  Then (11) becomes exactly the bound stated in Theorem~\ref{thm:main}:
\[
\bar{G}_T\le
\frac{2\big(f(x_0)-f^{\star}\big)}{\alpha T}
   +\frac{L\alpha d}{2}
   +\frac{2(1-\beta)G}{\sqrt{d}} .
\] 
\qedhere
\end{enumerate}

\section*{Rate Derivation (With Constants)}
Choosing $\alpha = \sqrt{\dfrac{2\big(f(x_0)-f^{\star}\big)}{L d\,T}}$ satisfies the stepsize restriction
$\alpha\le 1/(L\sqrt{d})$ (Assumption~\ref{ass:stepsize}) and yields
\[
\bar{G}_T
   \le
   2\sqrt{\frac{L d\big(f(x_0)-f^{\star}\big)}{T}}
   \;+\;
   \frac{2(1-\beta)G}{\sqrt{d}}
   \;=\;
   O\!\left(\frac{1}{\sqrt{T}}\right) + O\!\left(\frac{1-\beta}{\sqrt{d}}\right).
\]
If $\beta$ is taken close to $1$ (e.g. $\beta=0.9$ or $0.99$) the second term becomes negligible, and the overall rate matches that of full‑precision SGD.

\section*{Special Cases}
\begin{enumerate}[leftmargin=*]
\item \textbf{No momentum ($\beta=0$).}  
The bound reduces to
\[
\bar{G}_T\le\frac{2\Delta_0}{\alpha T}
            +\frac{L\alpha d}{2}
            +\frac{2G}{\sqrt{d}} .
\]
The additional $2G/\sqrt{d}$ term reflects the bias caused by directly signing noisy gradients.

\item \textbf{Deterministic gradients ($\sigma=0$).}  
When $g_t=\nabla f(x_t)$ exactly, Lemma~\ref{lem:momentum-bias} gives zero bias, and the bound simplifies to
\[
\bar{G}_T\le\frac{2\Delta_0}{\alpha T}
            +\frac{L\alpha d}{2},
\]
identical to the classic analysis of deterministic SignSGD with momentum.

\item \textbf{High‑dimensional regime.}  
If $d\gg T$, the term $L\alpha d/2$ dominates unless $\alpha$ is scaled as $1/(L\sqrt{d})$, which is precisely the restriction imposed by Assumption~\ref{ass:stepsize}.
\end{enumerate}

\end{document}